Our results support an architectural view of equivariance in variational quantum learning. Meyer et al. showed that incorporating task symmetry into a VQLM can improve generalization. We take this as a starting point and ask what remains to be designed once that symmetry has been imposed. The answer is that equivariance is not a complete ansatz specification: it constrains parameter sharing and responses to symmetry-related inputs, but not which trainable interactions the circuit should contain.

The subgroup experiments do not show that rotation-only equivariance is generally preferable to full board equivariance. They show that symmetry strength is a meaningful modeling choice. In this benchmark, weaker symmetry keeps most of the benefit of the fully symmetric model while giving the circuit slightly more freedom. Subgroup comparisons can therefore diagnose overconstraint when a fully equivariant circuit underperforms for a given data regime, circuit family, or optimizer.

The stronger empirical effect comes from changing the interactions inside the equivariant ansatz. The original edge circuit couples neighboring fields and must represent winning triples indirectly. The line interactions act directly on the three-cell structures that determine the label, while still sharing parameters across symmetry-related instances. This gives the largest improvement in our experiments. The matched random-sharing control rules out parameter reduction as the explanation: what matters is that trainable degrees of freedom are placed according to the task geometry.

This suggests a practical design principle. Once a task symmetry has been identified, one should ask not only whether the circuit respects it, but whether the equivariant gates expose the structures through which the label is determined. Equivariance is a starting point for ansatz design rather than its end point: it defines a constrained design space in which the choice of interactions can still be decisive.

Tic-Tac-Toe is small, enumerable, and exactly solvable by a deterministic rule oracle, so its value is as a controlled setting for isolating symmetry-aware architectural choices. The supported claim is narrow: task-aligned equivariant interactions can improve generalization while preserving the imposed symmetry. For larger problems, the transferable recipe is to identify task-relevant motifs and encode them as trainable interactions shared over the corresponding symmetry orbits.

A natural next step is to compare enriched equivariant ansatze with soft symmetry approaches, where controlled symmetry-breaking terms interpolate between hard equivariance and a free ansatz. Such comparisons should include not only the original symmetric circuit, but also equivariant circuits whose interactions have been designed to match the task. The broader lesson is that symmetry-aware quantum circuit design should use symmetry to organize the trainable interactions through which the model learns.
